% ===== PRÉAMBULE CANONIQUE — à coller VERBATIM en tête de CHAQUE .tex =====
\documentclass[11pt,a4paper]{article}
\usepackage[utf8]{inputenc}\usepackage[T1]{fontenc}\usepackage{lmodern}
\usepackage[french]{babel}
\usepackage{amsmath,amssymb,mathtools}\usepackage{icomma}
\usepackage[version=4]{mhchem}          % \ce{...} : équations & formules
\usepackage{siunitx}\sisetup{locale=FR} % \qty{}{}, \si{}, \ang{}
\usepackage{chemfig}                    % \chemfig{...} : structures développées
\usepackage{xcolor}\usepackage{colortbl}
\usepackage{tikz}\usepackage{pgfplots}\pgfplotsset{compat=1.18}
\usepackage[most]{tcolorbox}\usepackage{enumitem}\usepackage{array,booktabs}
\usepackage[a4paper,margin=2.2cm]{geometry}
\usetikzlibrary{arrows.meta,calc,angles,quotes,positioning,patterns,backgrounds,shapes.geometric}
\definecolor{primary}{RGB}{222,49,99}\definecolor{primarydark}{RGB}{150,22,55}
\definecolor{defcol}{RGB}{0,114,178}\definecolor{methcol}{RGB}{52,92,148}
\definecolor{excol}{RGB}{0,135,90}\definecolor{attncol}{RGB}{200,40,55}
\tcbset{boxbase/.style={enhanced,breakable,boxrule=0.9pt,arc=2.5pt,
  left=3mm,right=3mm,top=2.3mm,bottom=2.3mm,fonttitle=\bfseries,coltitle=white}}
% --- boîtes TUTORIEL (noms EXACTS, ne pas renommer) ---
\newtcolorbox{cle}[1][]{boxbase,colback=defcol!6,colframe=defcol,
  title={\textsc{À retenir}\ifx\relax#1\relax\relax\else\textmd{ : \itshape #1}\fi}}
\newtcolorbox{methode}[1][]{boxbase,colback=methcol!7,colframe=methcol,sharp corners,
  title={\textsc{Méthode}\ifx\relax#1\relax\relax\else\textmd{ --- \itshape #1}\fi}}
\newtcolorbox{exemplebox}[1][]{boxbase,colback=excol!5,colframe=excol,
  title={\textsc{Exemple}\ifx\relax#1\relax\relax\else\textmd{ : \itshape #1}\fi}}
\newtcolorbox{piege}[1][]{boxbase,colback=attncol!6,colframe=attncol,
  title={\textsc{Piège}\ifx\relax#1\relax\relax\else\textmd{ : \itshape #1}\fi}}
\newtcolorbox{remarque}[1][]{boxbase,colback=black!4,colframe=black!45,
  title={\textsc{Remarque}\ifx\relax#1\relax\relax\else\textmd{ : \itshape #1}\fi}}
% --- boîtes EXERCICE / CORRIGÉ (noms EXACTS, auto-comptées) ---
\newtcolorbox[auto counter]{exo}[1][]{boxbase,colback=primary!4,colframe=primary,
  title={\textsc{Exercice}~\thetcbcounter\ifx\relax#1\relax\relax\else\textmd{ --- \itshape #1}\fi}}
\newtcolorbox[auto counter]{corrige}[1][]{boxbase,colback=excol!4,colframe=excol,
  title={\textsc{Corrigé}~\thetcbcounter\ifx\relax#1\relax\relax\else\textmd{ --- \itshape #1}\fi}}
\newcommand{\R}{\mathbb{R}}\newcommand{\N}{\mathbb{N}}\newcommand{\Z}{\mathbb{Z}}\newcommand{\ds}{\displaystyle}
% (un \usepackage{fancyhdr}+en-tête est possible ; il est ignoré côté web)
% =========================================================================
\begin{document}
\sisetup{per-mode=symbol}
\begin{exo}[rendement d'une fermentation]
La fermentation du glucose s'écrit $\ce{C6H12O6 -> 2 C2H5OH + 2 CO2}$. On fait fermenter
\qty{0.5}{\mole} de glucose et l'on recueille \qty{11.5}{\gram} d'éthanol
($M=\qty{46}{\gram\per\mole}$).
\begin{enumerate}[label=\alph*)]
  \item Quelle quantité \emph{théorique} d'éthanol prévoit-on ?
  \item Quelle quantité d'éthanol obtient-on réellement ?
  \item En déduire le rendement de la fermentation.
\end{enumerate}
\end{exo}

\begin{exo}[calcination du calcaire]
On calcine \qty{100}{\gram} de calcaire \ce{CaCO3} : $\ce{CaCO3 -> CaO + CO2 ^}$. On donne
$M(\ce{CaCO3})=\qty{100}{\gram\per\mole}$, $M(\ce{CaO})=\qty{56}{\gram\per\mole}$.
\begin{enumerate}[label=\alph*)]
  \item Quelle masse \emph{théorique} de chaux \ce{CaO} peut-on obtenir ?
  \item On recueille \qty{42}{\gram} de \ce{CaO}. Quel est le rendement ?
\end{enumerate}
\end{exo}

\begin{exo}[remonter du produit voulu au réactif]
On veut réellement produire \qty{0.9}{\mole} d'ammoniac par $\ce{N2 + 3 H2 -> 2 NH3}$ (\ce{H2} en
excès), avec un rendement de \qty{75}{\percent}.
\begin{enumerate}[label=\alph*)]
  \item Quelle quantité \emph{théorique} de \ce{NH3} faut-il viser ?
  \item Quelle quantité de \ce{N2} faut-il alors engager ?
\end{enumerate}
\end{exo}

\begin{exo}[rendement à partir du limitant]
On plonge \qty{2.7}{\gram} d'aluminium ($M=\qty{27}{\gram\per\mole}$) dans \qty{200}{\milli\litre} de
\ce{ZnCl2} à \qty{0.5}{\mole\per\litre} : $\ce{3 ZnCl2 + 2 Al -> 3 Zn + 2 AlCl3}$. On donne
$M(\ce{Zn})=\qty{65.4}{\gram\per\mole}$.
\begin{enumerate}[label=\alph*)]
  \item Identifie le réactif limitant et la quantité \emph{théorique} de \ce{Zn}.
  \item Quelle masse \emph{théorique} de zinc cela représente-t-il ?
  \item On récupère \qty{3.27}{\gram} de \ce{Zn}. Quel est le rendement ?
\end{enumerate}
\end{exo}

\begin{exo}[synthèse du trioxyde de soufre]
Pour $\ce{2 SO2 + O2 -> 2 SO3}$, on engage \qty{128}{\gram} de \ce{SO2} dans un excès d'oxygène.
On donne $M(\ce{SO2})=\qty{64}{\gram\per\mole}$, $M(\ce{SO3})=\qty{80}{\gram\per\mole}$.
\begin{enumerate}[label=\alph*)]
  \item Quelle masse \emph{théorique} de \ce{SO3} peut-on obtenir ?
  \item On en recueille \qty{120}{\gram}. Quel est le rendement ?
\end{enumerate}
\end{exo}

\end{document}
